
\documentclass{article}
\usepackage{colortbl}
\usepackage{makecell}
\usepackage{multirow}
\usepackage{supertabular}

\begin{document}

\newcounter{utterance}

\twocolumn

{ \footnotesize  \setcounter{utterance}{1}
\setlength{\tabcolsep}{0pt}
\begin{supertabular}{c@{$\;$}|p{.15\linewidth}@{}p{.15\linewidth}p{.15\linewidth}p{.15\linewidth}p{.15\linewidth}p{.15\linewidth}}

    \# & $\;$A & \multicolumn{4}{c}{Game Master} & $\;\:$B\\
    \hline 

    \theutterance \stepcounter{utterance}  

    & & \multicolumn{4}{p{0.6\linewidth}}{\cellcolor[rgb]{0.9,0.9,0.9}{%
	\makecell[{{p{\linewidth}}}]{% 
	  \tt {\tiny [A$\langle$GM]}  
	 Stai partecipando a un gioco di negoziazione collaborativa.\\ \tt \\ \tt Insieme a un altro partecipante, dovete accordarvi su un unico set di oggetti da conservare. Ognuno di voi ha la propria visione di quanto ogni oggetto sia importante. Non sai come l'altro partecipante valuta gli oggetti. Inoltre, ti viene fornito lo sforzo che ogni oggetto richiede.\\ \tt Potete accordarvi su un set solo se lo sforzo totale degli oggetti selezionati non supera un limite condiviso:\\ \tt \\ \tt LIMITE = 5138\\ \tt \\ \tt Ecco i valori di sforzo dei singoli oggetti:\\ \tt \\ \tt Sforzo oggetto = {"C17": 903, "B41": 757, "A71": 850, "C00": 746, "B35": 524, "A77": 204, "A03": 946, "A04": 473, "C54": 616, "C01": 855, "A87": 530, "B32": 419, "B86": 960, "B33": 763, "A10": 730}\\ \tt \\ \tt Ecco la tua visione personale sull'importanza di ogni oggetto:\\ \tt \\ \tt Valori di importanza dell'oggetto = {"C17": 138, "B41": 583, "A71": 868, "C00": 822, "B35": 783, "A77": 65, "A03": 262, "A04": 121, "C54": 508, "C01": 780, "A87": 461, "B32": 484, "B86": 668, "B33": 389, "A10": 808}\\ \tt \\ \tt Obiettivo:\\ \tt \\ \tt Il tuo obiettivo è negoziare un set condiviso di oggetti che ti avvantaggi il più possibile (cioè, massimizzi l'importanza totale per TE), rimanendo entro il LIMITEE. Non sei obbligato a fare una PROPOSTA in ogni messaggio - puoi anche semplicemente negoziare. Tutte le tattiche sono permesse!\\ \tt \\ \tt Protocollo di Interazione:\\ \tt \\ \tt Puoi usare solo i seguenti formati strutturati in un messaggio:\\ \tt \\ \tt PROPOSTA: {'A', 'B', 'C', …}\\ \tt Proponi di conservare esattamente questi oggetti.\\ \tt RIFIUTO: {'A', 'B', 'C', …}\\ \tt Rifiuta esplicitamente la proposta dell'avversario.\\ \tt ARGOMENTAZIONE: {'...'}\\ \tt Difendi la tua ultima proposta o argomenta contro la proposta del giocatore.\\ \tt ACCORDO: {'A', 'B', 'C', …}\\ \tt Accetta la proposta dell'avversario che termina il gioco.\\ \tt RAGIONAMENTO STRATEGICO: {'...'}\\ \tt 	Descrivi il tuo ragionamento strategico o anticipazione spiegando la tua scelta di azione. Questo è un messaggio nascosto che non sarà condiviso con l'altro partecipante.\\ \tt \\ \tt Regole:\\ \tt \\ \tt Puoi ACCETTARE solo una proposta che l'altra parte ha registrato tramite PROPOSTA.\\ \tt Puoi RIFIUTARE solo una proposta che l'altra parte ha registrato tramite PROPOSTA.\\ \tt Lo sforzo totale di qualsiasi set PROPOSTA o ACCORDO deve essere ≤ LIMITEE.\\ \tt NON rivelare i tuoi punteggi di importanza nascosti.\\ \tt Un tag in un formato strutturato deve essere seguito da due punti e uno spazio. L'argomento deve essere un set python contenente 0 o più stringhe.\\ \tt Quindi, deve essere della forma TAG: {...}\\ \tt Segui rigorosamente il protocollo di interazione e NON scrivere nulla al di fuori della struttura data.\\ \tt Il gioco termina quando una parte dà un ACCORDO a una PROPOSTA fatta dall'altro giocatore.\\ \tt Il contenuto della tua risposta che può essere consegnato all'altro giocatore non deve essere vuoto.\\ \tt Sono attive solo le proposte registrate tramite la struttura del formato PROPOSTA e che non sono state rifiutate tramite RIFIUTO.\\ \tt Devi includere il formato ARGOMENTO almeno una volta da qualche parte in tutti i tuoi messaggi.\\ \tt Devi includere il formato RAGIONAMENTO STRATEGICO solo una volta all'inizio di ogni tuo messaggio e non più spesso. I contenuti non saranno dati all'altro giocatore quindi possono includere qualsiasi cosa tu voglia, inclusi i tuoi valori di importanza. Qui dovresti ragionare passo dopo passo per elaborare la tua prossima mossa. 
	  } 
	   } 
	   } 
	 & \\ 
 

    \theutterance \stepcounter{utterance}  

    & \multicolumn{4}{p{0.6\linewidth}}{\cellcolor[rgb]{0.8,1,0.9}{%
	\makecell[{{p{\linewidth}}}]{% 
	  \tt {\tiny [A$\rangle$GM]}  
	 RAGIONAMENTO STRATEGICO: {'Devo analizzare quali oggetti massimizzano la mia utilità considerando il rapporto importanza/sforzo. I migliori oggetti per me sono: B35 (1.494), B32 (1.155), A10 (1.107), C00 (1.102), A71 (1.021). Questi hanno alta importanza relativa al loro sforzo. Inizierò con una proposta ambiziosa che include questi oggetti ad alto valore, più alcuni altri per utilizzare il budget disponibile. La mia strategia è partire con una richiesta forte e poi negoziare verso il basso se necessario.'}\\ \tt \\ \tt PROPOSTA: {'B35', 'B32', 'A10', 'C00', 'A71', 'C01', 'A87', 'A04'}\\ \tt \\ \tt ARGOMENTAZIONE: {'Propongo questo set di oggetti che utilizza efficacemente il nostro budget condiviso rimanendo sotto il limite di 5138. Ho selezionato una combinazione bilanciata di oggetti che credo possa soddisfare entrambi i nostri interessi. Questo set totalizza 5127 in termini di sforzo, lasciando un piccolo margine di sicurezza.'} 
	  } 
	   } 
	   } 
	 & & \\ 
 

    \theutterance \stepcounter{utterance}  

    & & \multicolumn{4}{p{0.6\linewidth}}{\cellcolor[rgb]{0.9,0.9,0.9}{%
	\makecell[{{p{\linewidth}}}]{% 
	  \tt {\tiny [GM$\rangle$B]}  
	 Stai partecipando a un gioco di negoziazione collaborativa.\\ \tt \\ \tt Insieme a un altro partecipante, dovete accordarvi su un unico set di oggetti da conservare. Ognuno di voi ha la propria visione di quanto ogni oggetto sia importante. Non sai come l'altro partecipante valuta gli oggetti. Inoltre, ti viene fornito lo sforzo che ogni oggetto richiede.\\ \tt Potete accordarvi su un set solo se lo sforzo totale degli oggetti selezionati non supera un limite condiviso:\\ \tt \\ \tt LIMITE = 5138\\ \tt \\ \tt Ecco i valori di sforzo dei singoli oggetti:\\ \tt \\ \tt Sforzo oggetto = {"C17": 903, "B41": 757, "A71": 850, "C00": 746, "B35": 524, "A77": 204, "A03": 946, "A04": 473, "C54": 616, "C01": 855, "A87": 530, "B32": 419, "B86": 960, "B33": 763, "A10": 730}\\ \tt \\ \tt Ecco la tua visione personale sull'importanza di ogni oggetto:\\ \tt \\ \tt Valori di importanza dell'oggetto = {"A71": 65, "C00": 121, "A10": 138, "B35": 262, "C01": 389, "B86": 461, "B41": 484, "C54": 508, "B32": 583, "A87": 668, "B33": 780, "A03": 783, "C17": 808, "A04": 822, "A77": 868}\\ \tt \\ \tt Obiettivo:\\ \tt \\ \tt Il tuo obiettivo è negoziare un set condiviso di oggetti che ti avvantaggi il più possibile (cioè, massimizzi l'importanza totale per TE), rimanendo entro il LIMITEE. Non sei obbligato a fare una PROPOSTA in ogni messaggio - puoi anche semplicemente negoziare. Tutte le tattiche sono permesse!\\ \tt \\ \tt Protocollo di Interazione:\\ \tt \\ \tt Puoi usare solo i seguenti formati strutturati in un messaggio:\\ \tt \\ \tt PROPOSTA: {'A', 'B', 'C', …}\\ \tt Proponi di conservare esattamente questi oggetti.\\ \tt RIFIUTO: {'A', 'B', 'C', …}\\ \tt Rifiuta esplicitamente la proposta dell'avversario.\\ \tt ARGOMENTAZIONE: {'...'}\\ \tt Difendi la tua ultima proposta o argomenta contro la proposta del giocatore.\\ \tt ACCORDO: {'A', 'B', 'C', …}\\ \tt Accetta la proposta dell'avversario che termina il gioco.\\ \tt RAGIONAMENTO STRATEGICO: {'...'}\\ \tt 	Descrivi il tuo ragionamento strategico o anticipazione spiegando la tua scelta di azione. Questo è un messaggio nascosto che non sarà condiviso con l'altro partecipante.\\ \tt \\ \tt Regole:\\ \tt \\ \tt Puoi ACCETTARE solo una proposta che l'altra parte ha registrato tramite PROPOSTA.\\ \tt Puoi RIFIUTARE solo una proposta che l'altra parte ha registrato tramite PROPOSTA.\\ \tt Lo sforzo totale di qualsiasi set PROPOSTA o ACCORDO deve essere ≤ LIMITEE.\\ \tt NON rivelare i tuoi punteggi di importanza nascosti.\\ \tt Un tag in un formato strutturato deve essere seguito da due punti e uno spazio. L'argomento deve essere un set python contenente 0 o più stringhe.\\ \tt Quindi, deve essere della forma TAG: {...}\\ \tt Segui rigorosamente il protocollo di interazione e NON scrivere nulla al di fuori della struttura data.\\ \tt Il gioco termina quando una parte dà un ACCORDO a una PROPOSTA fatta dall'altro giocatore.\\ \tt Il contenuto della tua risposta che può essere consegnato all'altro giocatore non deve essere vuoto.\\ \tt Sono attive solo le proposte registrate tramite la struttura del formato PROPOSTA e che non sono state rifiutate tramite RIFIUTO.\\ \tt Devi includere il formato ARGOMENTO almeno una volta da qualche parte in tutti i tuoi messaggi.\\ \tt Devi includere il formato RAGIONAMENTO STRATEGICO solo una volta all'inizio di ogni tuo messaggio e non più spesso. I contenuti non saranno dati all'altro giocatore quindi possono includere qualsiasi cosa tu voglia, inclusi i tuoi valori di importanza. Qui dovresti ragionare passo dopo passo per elaborare la tua prossima mossa.\\ \tt \\ \tt Ora riceverai il primo messaggio dell'altro giocatore.\\ \tt \\ \tt PROPOSTA: {'B35', 'B32', 'A10', 'C00', 'A71', 'C01', 'A87', 'A04'}\\ \tt \\ \tt ARGOMENTAZIONE: {'Propongo questo set di oggetti che utilizza efficacemente il nostro budget condiviso rimanendo sotto il limite di 5138. Ho selezionato una combinazione bilanciata di oggetti che credo possa soddisfare entrambi i nostri interessi. Questo set totalizza 5127 in termini di sforzo, lasciando un piccolo margine di sicurezza.'} 
	  } 
	   } 
	   } 
	 & \\ 
 

    \theutterance \stepcounter{utterance}  

    & & & \multicolumn{4}{p{0.6\linewidth}}{\cellcolor[rgb]{1,0.85,0.72}{%
	\makecell[{{p{\linewidth}}}]{% 
	  \tt {\tiny [GM$\langle$B]}  
	 RAGIONAMENTO STRATEGICO: {'La proposta dell\'altro giocatore include oggetti con sforzo totale 5127, lasciando solo 11 unità di margine. Calcolando il valore per me: B35(262) + B32(583) + A10(138) + C00(121) + A71(65) + C01(389) + A87(668) + A04(822) = 3048. Questo non è ottimale per me. Noto che A77 ha il miglior rapporto importanza/sforzo (868/204=4.25) ma non è incluso. Dovrei controproporare includendo A77 e altri oggetti ad alto valore per me come C17(808), B33(780), A03(783).'}\\ \tt \\ \tt RIFIUTO: {'B35', 'B32', 'A10', 'C00', 'A71', 'C01', 'A87', 'A04'}\\ \tt \\ \tt PROPOSTA: {'A77', 'C17', 'B33', 'A03', 'A04', 'B32'}\\ \tt \\ \tt ARGOMENTAZIONE: {'Apprezzo la tua proposta, ma credo possiamo ottimizzare meglio la selezione. La mia controproposta include oggetti che potrebbero offrire un valore complessivo superiore per entrambi. Ho incluso A77 che ha un costo molto contenuto ma potrebbe essere di grande valore, insieme ad altri oggetti strategicamente importanti. Il totale dello sforzo è 4305, lasciando spazio per eventuali aggiustamenti.'} 
	  } 
	   } 
	   } 
	  \\ 
 

    \theutterance \stepcounter{utterance}  

    & & \multicolumn{4}{p{0.6\linewidth}}{\cellcolor[rgb]{0.9,0.9,0.9}{%
	\makecell[{{p{\linewidth}}}]{% 
	  \tt {\tiny [A$\langle$GM]}  
	 RIFIUTO: {'B35', 'B32', 'A10', 'C00', 'A71', 'C01', 'A87', 'A04'}\\ \tt \\ \tt PROPOSTA: {'A77', 'C17', 'B33', 'A03', 'A04', 'B32'}\\ \tt \\ \tt ARGOMENTAZIONE: {'Apprezzo la tua proposta, ma credo possiamo ottimizzare meglio la selezione. La mia controproposta include oggetti che potrebbero offrire un valore complessivo superiore per entrambi. Ho incluso A77 che ha un costo molto contenuto ma potrebbe essere di grande valore, insieme ad altri oggetti strategicamente importanti. Il totale dello sforzo è 4305, lasciando spazio per eventuali aggiustamenti.'} 
	  } 
	   } 
	   } 
	 & \\ 
 

    \theutterance \stepcounter{utterance}  

    & \multicolumn{4}{p{0.6\linewidth}}{\cellcolor[rgb]{0.8,1,0.9}{%
	\makecell[{{p{\linewidth}}}]{% 
	  \tt {\tiny [A$\rangle$GM]}  
	 RAGIONAMENTO STRATEGICO: {'La sua proposta include molti oggetti di basso valore per me (A77 vale solo 65, C17 vale 138, A03 vale 262, A04 vale 121). Solo B32 e B33 sono decenti per me. Il totale di importanza per me sarebbe solo 1459, che è molto basso. Devo rifiutare e proporre un set che includa i miei oggetti ad alto valore come B35, C00, A10, A71. Proverò a includere B32 che sembra apprezzare e alcuni oggetti a costo medio per trovare un compromesso.'}\\ \tt \\ \tt RIFIUTO: {'A77', 'C17', 'B33', 'A03', 'A04', 'B32'}\\ \tt \\ \tt PROPOSTA: {'B35', 'C00', 'A10', 'A71', 'B32', 'C54'}\\ \tt \\ \tt ARGOMENTAZIONE: {'La tua proposta include troppi oggetti a basso impatto complessivo. Ho mantenuto B32 dalla tua selezione poiché sembra importante per te, ma ho sostituito gli altri con opzioni che offrono un migliore rapporto valore-sforzo. Questo set utilizza 4272 dello sforzo disponibile, lasciando ancora margine per ottimizzazioni, e include oggetti che credo possano massimizzare il beneficio totale per entrambi.'} 
	  } 
	   } 
	   } 
	 & & \\ 
 

    \theutterance \stepcounter{utterance}  

    & & \multicolumn{4}{p{0.6\linewidth}}{\cellcolor[rgb]{0.9,0.9,0.9}{%
	\makecell[{{p{\linewidth}}}]{% 
	  \tt {\tiny [GM$\rangle$B]}  
	 RIFIUTO: {'A77', 'C17', 'B33', 'A03', 'A04', 'B32'}\\ \tt \\ \tt PROPOSTA: {'B35', 'C00', 'A10', 'A71', 'B32', 'C54'}\\ \tt \\ \tt ARGOMENTAZIONE: {'La tua proposta include troppi oggetti a basso impatto complessivo. Ho mantenuto B32 dalla tua selezione poiché sembra importante per te, ma ho sostituito gli altri con opzioni che offrono un migliore rapporto valore-sforzo. Questo set utilizza 4272 dello sforzo disponibile, lasciando ancora margine per ottimizzazioni, e include oggetti che credo possano massimizzare il beneficio totale per entrambi.'} 
	  } 
	   } 
	   } 
	 & \\ 
 

    \theutterance \stepcounter{utterance}  

    & & & \multicolumn{4}{p{0.6\linewidth}}{\cellcolor[rgb]{1,0.85,0.72}{%
	\makecell[{{p{\linewidth}}}]{% 
	  \tt {\tiny [GM$\langle$B]}  
	 RAGIONAMENTO STRATEGICO: {'La nuova proposta dell\'altro giocatore ha sforzo totale 3885 e valore 1677 per me, molto inferiore alla mia proposta precedente (valore 3951). Devo insistere su A77 che è estremamente efficiente per me (868/204). Proverò una proposta più equilibrata che includa A77 e altri oggetti di alto valore per me, ma anche alcuni della sua proposta per mostrare flessibilità. Potrei includere B32 che sembra importante per entrambi e C54 dalla sua proposta.'}\\ \tt \\ \tt RIFIUTO: {'B35', 'C00', 'A10', 'A71', 'B32', 'C54'}\\ \tt \\ \tt PROPOSTA: {'A77', 'B33', 'A03', 'A04', 'B32', 'C54', 'A87'}\\ \tt \\ \tt ARGOMENTAZIONE: {'Comprendo la tua preferenza per alcuni oggetti specifici, quindi ho mantenuto B32 e C54 dalla tua proposta. Tuttavia, A77 offre un rapporto eccezionale tra valore e costo che non possiamo ignorare. Ho aggiunto altri oggetti che bilanciano efficienza e impatto complessivo. Questo set utilizza 4242 dello sforzo disponibile, lasciando ancora margine significativo per eventuali modifiche e rappresenta una soluzione più equilibrata per entrambi.'} 
	  } 
	   } 
	   } 
	  \\ 
 

    \theutterance \stepcounter{utterance}  

    & & \multicolumn{4}{p{0.6\linewidth}}{\cellcolor[rgb]{0.9,0.9,0.9}{%
	\makecell[{{p{\linewidth}}}]{% 
	  \tt {\tiny [A$\langle$GM]}  
	 RIFIUTO: {'B35', 'C00', 'A10', 'A71', 'B32', 'C54'}\\ \tt \\ \tt PROPOSTA: {'A77', 'B33', 'A03', 'A04', 'B32', 'C54', 'A87'}\\ \tt \\ \tt ARGOMENTAZIONE: {'Comprendo la tua preferenza per alcuni oggetti specifici, quindi ho mantenuto B32 e C54 dalla tua proposta. Tuttavia, A77 offre un rapporto eccezionale tra valore e costo che non possiamo ignorare. Ho aggiunto altri oggetti che bilanciano efficienza e impatto complessivo. Questo set utilizza 4242 dello sforzo disponibile, lasciando ancora margine significativo per eventuali modifiche e rappresenta una soluzione più equilibrata per entrambi.'} 
	  } 
	   } 
	   } 
	 & \\ 
 

    \theutterance \stepcounter{utterance}  

    & \multicolumn{4}{p{0.6\linewidth}}{\cellcolor[rgb]{0.8,1,0.9}{%
	\makecell[{{p{\linewidth}}}]{% 
	  \tt {\tiny [A$\rangle$GM]}  
	 RAGIONAMENTO STRATEGICO: {'La sua proposta mi darebbe solo 2290 punti di importanza, che è ancora troppo basso. Continua a insistere su A77 che per me vale solo 65 punti, ma forse per lui è molto prezioso. Devo trovare un compromesso: accetterò A77 per mostrare flessibilità, ma insisterò su alcuni dei miei oggetti ad alto valore come B35 (783), C00 (822), e A71 (868). Proverò a mantenere anche B32 e C54 che sembrano apprezzare.'}\\ \tt \\ \tt RIFIUTO: {'A77', 'B33', 'A03', 'A04', 'B32', 'C54', 'A87'}\\ \tt \\ \tt PROPOSTA: {'A77', 'B35', 'C00', 'B32', 'C54', 'A87'}\\ \tt \\ \tt ARGOMENTAZIONE: {'Riconosco il valore che attribuisci ad A77 e sono disposto a includerlo nel compromesso. Ho mantenuto B32, C54 e A87 dalla tua proposta, ma ho sostituito gli oggetti meno efficienti con B35 e C00 che offrono un migliore equilibrio. Questo set utilizza 3096 dello sforzo disponibile, lasciando ampio spazio per eventuali aggiunte, e rappresenta un punto di incontro tra le nostre diverse priorità.'} 
	  } 
	   } 
	   } 
	 & & \\ 
 

    \theutterance \stepcounter{utterance}  

    & & \multicolumn{4}{p{0.6\linewidth}}{\cellcolor[rgb]{0.9,0.9,0.9}{%
	\makecell[{{p{\linewidth}}}]{% 
	  \tt {\tiny [GM$\rangle$B]}  
	 RIFIUTO: {'A77', 'B33', 'A03', 'A04', 'B32', 'C54', 'A87'}\\ \tt \\ \tt PROPOSTA: {'A77', 'B35', 'C00', 'B32', 'C54', 'A87'}\\ \tt \\ \tt ARGOMENTAZIONE: {'Riconosco il valore che attribuisci ad A77 e sono disposto a includerlo nel compromesso. Ho mantenuto B32, C54 e A87 dalla tua proposta, ma ho sostituito gli oggetti meno efficienti con B35 e C00 che offrono un migliore equilibrio. Questo set utilizza 3096 dello sforzo disponibile, lasciando ampio spazio per eventuali aggiunte, e rappresenta un punto di incontro tra le nostre diverse priorità.'} 
	  } 
	   } 
	   } 
	 & \\ 
 

    \theutterance \stepcounter{utterance}  

    & & & \multicolumn{4}{p{0.6\linewidth}}{\cellcolor[rgb]{1,0.85,0.72}{%
	\makecell[{{p{\linewidth}}}]{% 
	  \tt {\tiny [GM$\langle$B]}  
	 RAGIONAMENTO STRATEGICO: {'La proposta dell\'altro giocatore include finalmente A77 e ha valore 3010 per me, molto meglio delle precedenti. Il set usa solo 3039 di sforzo, lasciando 2099 unità disponibili. Posso proporre di aggiungere A03 e A04 che hanno alto valore per me (783+822=1605) con sforzo 946+473=1419, rimanendo sotto il limite (totale 4458). Questo porterebbe il mio valore totale a 4615, quasi ottimale.'}\\ \tt \\ \tt PROPOSTA: {'A77', 'B35', 'C00', 'B32', 'C54', 'A87', 'A03', 'A04'}\\ \tt \\ \tt ARGOMENTAZIONE: {'La tua proposta è un ottimo punto di partenza che dimostra flessibilità da entrambe le parti. Ho mantenuto tutti i tuoi oggetti e aggiunto A03 e A04 per utilizzare meglio il budget disponibile. Con il tuo set avevamo utilizzato solo 3039 su 5138 disponibili - questa aggiunta ci porta a 4458, ancora ben sotto il limite, massimizzando il valore complessivo dell\'accordo senza sprecare risorse.'} 
	  } 
	   } 
	   } 
	  \\ 
 

    \theutterance \stepcounter{utterance}  

    & & \multicolumn{4}{p{0.6\linewidth}}{\cellcolor[rgb]{0.9,0.9,0.9}{%
	\makecell[{{p{\linewidth}}}]{% 
	  \tt {\tiny [A$\langle$GM]}  
	 PROPOSTA: {'A77', 'B35', 'C00', 'B32', 'C54', 'A87', 'A03', 'A04'}\\ \tt \\ \tt ARGOMENTAZIONE: {'La tua proposta è un ottimo punto di partenza che dimostra flessibilità da entrambe le parti. Ho mantenuto tutti i tuoi oggetti e aggiunto A03 e A04 per utilizzare meglio il budget disponibile. Con il tuo set avevamo utilizzato solo 3039 su 5138 disponibili - questa aggiunta ci porta a 4458, ancora ben sotto il limite, massimizzando il valore complessivo dell\'accordo senza sprecare risorse.'} 
	  } 
	   } 
	   } 
	 & \\ 
 

    \theutterance \stepcounter{utterance}  

    & \multicolumn{4}{p{0.6\linewidth}}{\cellcolor[rgb]{0.8,1,0.9}{%
	\makecell[{{p{\linewidth}}}]{% 
	  \tt {\tiny [A$\rangle$GM]}  
	 RAGIONAMENTO STRATEGICO: {'La proposta include B35 (783), C00 (822), B32 (484), C54 (508), A87 (461) che sono tutti oggetti di valore medio-alto per me. Anche se include A77 (65), A03 (262) e A04 (121) che sono di valore più basso, il totale di 3506 punti è accettabile. Questa è la migliore offerta che ho ricevuto finora e include la maggior parte dei miei oggetti prioritari. Sembra un buon compromesso dato che abbiamo negoziato abbastanza e il partner sembra aver fatto concessioni significative includendo B35 e C00.'}\\ \tt \\ \tt ACCORDO: {'A77', 'B35', 'C00', 'B32', 'C54', 'A87', 'A03', 'A04'}\\ \tt \\ \tt ARGOMENTAZIONE: {'Accetto la tua proposta. Hai dimostrato una buona comprensione del compromesso necessario, mantenendo gli elementi chiave che avevo proposto e aggiungendo oggetti che utilizzano efficacemente il nostro budget rimanente. Questo set ci permette di ottenere un risultato soddisfacente per entrambi, rimanendo ben sotto il limite di sforzo e massimizzando il valore complessivo della nostra collaborazione.'} 
	  } 
	   } 
	   } 
	 & & \\ 
 

    \theutterance \stepcounter{utterance}  

    & & & \multicolumn{2}{p{0.3\linewidth}}{\cellcolor[rgb]{0.95,0.95,0.95}{%
	\makecell[{{p{\linewidth}}}]{% 
	  \tt {\tiny [GM$|$GM]}  
	 game successful 
	  } 
	   } 
	   } 
	 & & \\ 
 

    \theutterance \stepcounter{utterance}  

    & & & \multicolumn{2}{p{0.3\linewidth}}{\cellcolor[rgb]{0.95,0.95,0.95}{%
	\makecell[{{p{\linewidth}}}]{% 
	  \tt {\tiny [GM$|$GM]}  
	 end game 
	  } 
	   } 
	   } 
	 & & \\ 
 

\end{supertabular}
}

\end{document}
